\par\medskip\noindent
\textbf{\textit{\SGARefLink{sga3:xix:lemma:3:5}{Lemma 3.5}.}}---
\textit{Let \(S\) be a scheme, let \(G\) be a reductive \(S\)-group
of semisimple rank one, and let \(T\) be a maximal torus of \(G\).  If
\(\alpha\) is a root of \(G\) with respect to \(T\), then so is
\(-\alpha\), and}
\[
  \mathfrak g
  =
  \mathfrak t
  \oplus \mathfrak g^\alpha
  \oplus \mathfrak g^{-\alpha},
\]
\textit{where \(\mathfrak g^\alpha\) and
\(\mathfrak g^{-\alpha}\) are locally free of rank one.}
\SGARefTarget{sga3:xix:lemma:3:5}

\par\medskip\noindent
\SGAInlineRefTarget{sga3:xix:definition:3-6}\textbf{\textit{\SGARefLink{sga3:xix:definition:3:6}{Definition 3.6}.}}---
\textit{Let \(S\) be a scheme, let \(G\) be a reductive \(S\)-group,
let \(T\) be a maximal torus of \(G\), and let \(R\) be a set of roots
of \(G\) with respect to \(T\).  We say that \(R\) is a \emph{root
system} of \(G\) with respect to \(T\) if the following equivalent
conditions are satisfied:}
\SGARefTarget{sga3:xix:definition:3:6}

\begin{enumerate}
\item[\textit{(i)}]
\textit{For every \(s\in S\), the map
\(\alpha\mapsto\alpha_s\) is a bijection from \(R\) onto the set of
roots of \(G_s\) with respect to \(T_s\).}

\item[\textit{\SGAInlineRefTarget{sga3:xix:definition:3:6:item:ii}(ii)}]
\textit{The elements of \(R\) are distinct on every fiber
\textup{(that is, if \(\alpha,\alpha'\in R\) and
\(\alpha\ne\alpha'\), then
\(\alpha-\alpha'=\alpha{\alpha'}^{-1}\) is nontrivial on every
fiber)}, and, for every \(s\in S\),}
\[
  \dim(G_s/\kappa(s))-\dim(T_s/\kappa(s))=\operatorname{Card}(R).
\]

\item[\textit{\SGAInlineRefTarget{sga3:xix:definition:3:6:item:iii}(iii)}]
\textit{One has}
\[
  \mathfrak g
  =
  \mathfrak t
  \oplus
  \bigoplus_{\alpha\in R}\mathfrak g^\alpha.
\]
\end{enumerate}

The equivalence of these conditions is immediate.

\par\medskip\noindent
\SGAInlineRefTarget{sga3:xix:lemma:3-7}\textbf{\textit{\SGARefLink{sga3:xix:lemma:3:7}{Lemma 3.7}.}}---
\textit{Let \(S\) be a scheme, let \(G\) be a reductive \(S\)-group,
let \(T\) be a maximal torus of \(G\), and let \(R\) be a root system
of \(G\) with respect to \(T\).  Every root of \(G\) with respect to
\(T\) is locally on \(S\) equal to an element of \(R\).}
\SGARefTarget{sga3:xix:lemma:3:7}

\SGAProof{Proof}
This follows from condition \SGARefLink{sga3:xix:definition:3:6:item:iii}{(iii)}.

Put
\[
  \mathcal M
  =
  \underline{\operatorname{Hom}}_{S\text{-gr.}}
  (T,\mathbb G_{m,S});
\]
this is a twisted constant \(S\)-group scheme (\SGARefLink{sga3:x:theorem:5-6}{Exposé X, 5.6}).  If
\(G\) has a root system \(R\) with respect to \(T\), then the inclusion
\(R\hookrightarrow\mathcal M(S)\) defines a morphism
\[
  R_S\longrightarrow\mathcal M,
\]
where \(R_S\) is the constant \(S\)-scheme defined by \(R\).  By
condition \SGARefLink{sga3:xix:definition:3:6:item:ii}{(ii)}, this morphism is an open and closed immersion whose
image is precisely
\[
  \bigcup_{\alpha\in R}\alpha(S),
\]
where each \(\alpha\in R\) is regarded as a section
\(S\to\mathcal M\).

Let \(\mathcal R\) be the functor of roots of \(G\) with respect to
\(T\).  By definition, \(\mathcal R(S')\) is the set of roots of
\(G_{S'}\) with respect to \(T_{S'}\), for every \(S'\to S\).  If
\(S'=\varnothing\), set \(\mathcal R(\varnothing)=\{e\}\); if
\(S'\ne\varnothing\), the inclusion
\(R\hookrightarrow\mathcal M(S')\) identifies \(R\) with a root
system of \(G_{S'}\) with respect to \(T_{S'}\).  Thus, by \SGARefLink{sga3:xix:lemma:3-7}{3.7},
\[
  \mathcal R(S')
  =
  \operatorname{Hom}_{\mathrm{loc.\ const.}}(S',R),
\]
which shows that \(\mathcal R\) is represented by \(R_S\).

Now drop the assumption that \(G\) necessarily has a root system with
respect to \(T\).  The functor \(\mathcal R\) is in any case an fpqc
subsheaf of \(\mathcal M\).  Locally for this topology, \(G\) does
have a root system with respect to \(T\) (for example, make \(T\)
split and use the proof of \SGARefLink{sga3:xix:theorem:2:5}{Theorem 2.5}).  By \SGARefLink{sga3:iv:proposition:4:6:8}{Exposé IV, 4.6.8} and the
descent theory of open (respectively closed)
subschemes, we obtain the following proposition.
